\documentclass[11pt,a4paper]{article}
% Preambulo comun de los reportes de Tecnologias Celulares.
% Cada main.tex hace: \documentclass ... % Preambulo comun de los reportes de Tecnologias Celulares.
% Cada main.tex hace: \documentclass ... % Preambulo comun de los reportes de Tecnologias Celulares.
% Cada main.tex hace: \documentclass ... \input{../comun/preambulo.tex}
\usepackage[utf8]{inputenc}
\usepackage[T1]{fontenc}
\usepackage[spanish,mexico]{babel}
\usepackage{graphicx}
\usepackage[margin=2.5cm]{geometry}
\usepackage{amsmath}
\usepackage{amssymb}
\usepackage{booktabs}
\usepackage{tabularx}
\usepackage{array}
\usepackage{float}
\usepackage{setspace}
\usepackage{parskip}
\usepackage{fancyhdr}
\usepackage{textcomp}
\usepackage{xcolor}
\usepackage{caption}
\usepackage{fvextra}
\usepackage[hidelinks]{hyperref}

\onehalfspacing
\setlength{\parskip}{6pt}
\sloppy
\setlength{\emergencystretch}{3em}
\setlength{\headheight}{40pt}
\captionsetup{font=small,labelfont=bf}
\setlength{\tabcolsep}{6pt}

\pagestyle{fancy}
\fancyhf{}
\fancyhead[L]{\includegraphics[width=1cm]{/home/yovick/repos/UAZ-IEI/Assets/UIE.png}}
\fancyhead[R]{\small Tecnologías Celulares}
\fancyfoot[C]{\thepage}
\renewcommand{\headrulewidth}{0.4pt}

\DefineVerbatimEnvironment{codigo}{Verbatim}{%
  breaklines=true,breakanywhere=true,fontsize=\small,frame=single,
  framesep=3pt,rulecolor=\color{gray!50}}

% Figura centrada con pie: \figura[ancho]{archivo}{leyenda}
\newcommand{\figura}[3][0.85]{%
  \begin{center}
    \includegraphics[width=#1\linewidth]{#2}\\[2pt]
    {\small #3}
  \end{center}}

% Portada: \portada{numero}{titulo de la practica}
\newcommand{\portada}[2]{%
\begin{titlepage}
  \thispagestyle{empty}
  \begin{center}
    \vspace*{1cm}
    \includegraphics[width=4.5cm]{/home/yovick/repos/UAZ-IEI/Assets/uaz.png}
    \vspace{0.8cm}

    {\large Universidad Autónoma de Zacatecas}\\[4pt]
    {\large Unidad Académica de Ingeniería Eléctrica}\\[4pt]
    {\normalsize Ingeniería Electrónica Industrial con especialidad en Telecomunicaciones}

    \vspace{1.8cm}
    {\LARGE\bfseries Práctica #1}\\[10pt]
    {\Large\bfseries #2}
  \end{center}

  \vspace{2.2cm}
  \begin{center}
    \begin{tabular}{rl}
      \textbf{Materia:} & Tecnologías Celulares\\
      \textbf{Docente:} & Dr. Víktor Iván Rodríguez Abdalá\\
      \textbf{Alumno:}  & Arnold Yovick Ríos Zaldívar\\
      \textbf{Fecha:}   & Noviembre de 2026\\
    \end{tabular}
  \end{center}
  \vfill
\end{titlepage}}

\usepackage[utf8]{inputenc}
\usepackage[T1]{fontenc}
\usepackage[spanish,mexico]{babel}
\usepackage{graphicx}
\usepackage[margin=2.5cm]{geometry}
\usepackage{amsmath}
\usepackage{amssymb}
\usepackage{booktabs}
\usepackage{tabularx}
\usepackage{array}
\usepackage{float}
\usepackage{setspace}
\usepackage{parskip}
\usepackage{fancyhdr}
\usepackage{textcomp}
\usepackage{xcolor}
\usepackage{caption}
\usepackage{fvextra}
\usepackage[hidelinks]{hyperref}

\onehalfspacing
\setlength{\parskip}{6pt}
\sloppy
\setlength{\emergencystretch}{3em}
\setlength{\headheight}{40pt}
\captionsetup{font=small,labelfont=bf}
\setlength{\tabcolsep}{6pt}

\pagestyle{fancy}
\fancyhf{}
\fancyhead[L]{\includegraphics[width=1cm]{/home/yovick/repos/UAZ-IEI/Assets/UIE.png}}
\fancyhead[R]{\small Tecnologías Celulares}
\fancyfoot[C]{\thepage}
\renewcommand{\headrulewidth}{0.4pt}

\DefineVerbatimEnvironment{codigo}{Verbatim}{%
  breaklines=true,breakanywhere=true,fontsize=\small,frame=single,
  framesep=3pt,rulecolor=\color{gray!50}}

% Figura centrada con pie: \figura[ancho]{archivo}{leyenda}
\newcommand{\figura}[3][0.85]{%
  \begin{center}
    \includegraphics[width=#1\linewidth]{#2}\\[2pt]
    {\small #3}
  \end{center}}

% Portada: \portada{numero}{titulo de la practica}
\newcommand{\portada}[2]{%
\begin{titlepage}
  \thispagestyle{empty}
  \begin{center}
    \vspace*{1cm}
    \includegraphics[width=4.5cm]{/home/yovick/repos/UAZ-IEI/Assets/uaz.png}
    \vspace{0.8cm}

    {\large Universidad Autónoma de Zacatecas}\\[4pt]
    {\large Unidad Académica de Ingeniería Eléctrica}\\[4pt]
    {\normalsize Ingeniería Electrónica Industrial con especialidad en Telecomunicaciones}

    \vspace{1.8cm}
    {\LARGE\bfseries Práctica #1}\\[10pt]
    {\Large\bfseries #2}
  \end{center}

  \vspace{2.2cm}
  \begin{center}
    \begin{tabular}{rl}
      \textbf{Materia:} & Tecnologías Celulares\\
      \textbf{Docente:} & Dr. Víktor Iván Rodríguez Abdalá\\
      \textbf{Alumno:}  & Arnold Yovick Ríos Zaldívar\\
      \textbf{Fecha:}   & Noviembre de 2026\\
    \end{tabular}
  \end{center}
  \vfill
\end{titlepage}}

\usepackage[utf8]{inputenc}
\usepackage[T1]{fontenc}
\usepackage[spanish,mexico]{babel}
\usepackage{graphicx}
\usepackage[margin=2.5cm]{geometry}
\usepackage{amsmath}
\usepackage{amssymb}
\usepackage{booktabs}
\usepackage{tabularx}
\usepackage{array}
\usepackage{float}
\usepackage{setspace}
\usepackage{parskip}
\usepackage{fancyhdr}
\usepackage{textcomp}
\usepackage{xcolor}
\usepackage{caption}
\usepackage{fvextra}
\usepackage[hidelinks]{hyperref}

\onehalfspacing
\setlength{\parskip}{6pt}
\sloppy
\setlength{\emergencystretch}{3em}
\setlength{\headheight}{40pt}
\captionsetup{font=small,labelfont=bf}
\setlength{\tabcolsep}{6pt}

\pagestyle{fancy}
\fancyhf{}
\fancyhead[L]{\includegraphics[width=1cm]{/home/yovick/repos/UAZ-IEI/Assets/UIE.png}}
\fancyhead[R]{\small Tecnologías Celulares}
\fancyfoot[C]{\thepage}
\renewcommand{\headrulewidth}{0.4pt}

\DefineVerbatimEnvironment{codigo}{Verbatim}{%
  breaklines=true,breakanywhere=true,fontsize=\small,frame=single,
  framesep=3pt,rulecolor=\color{gray!50}}

% Figura centrada con pie: \figura[ancho]{archivo}{leyenda}
\newcommand{\figura}[3][0.85]{%
  \begin{center}
    \includegraphics[width=#1\linewidth]{#2}\\[2pt]
    {\small #3}
  \end{center}}

% Portada: \portada{numero}{titulo de la practica}
\newcommand{\portada}[2]{%
\begin{titlepage}
  \thispagestyle{empty}
  \begin{center}
    \vspace*{1cm}
    \includegraphics[width=4.5cm]{/home/yovick/repos/UAZ-IEI/Assets/uaz.png}
    \vspace{0.8cm}

    {\large Universidad Autónoma de Zacatecas}\\[4pt]
    {\large Unidad Académica de Ingeniería Eléctrica}\\[4pt]
    {\normalsize Ingeniería Electrónica Industrial con especialidad en Telecomunicaciones}

    \vspace{1.8cm}
    {\LARGE\bfseries Práctica #1}\\[10pt]
    {\Large\bfseries #2}
  \end{center}

  \vspace{2.2cm}
  \begin{center}
    \begin{tabular}{rl}
      \textbf{Materia:} & Tecnologías Celulares\\
      \textbf{Docente:} & Dr. Víktor Iván Rodríguez Abdalá\\
      \textbf{Alumno:}  & Arnold Yovick Ríos Zaldívar\\
      \textbf{Fecha:}   & Noviembre de 2026\\
    \end{tabular}
  \end{center}
  \vfill
\end{titlepage}}


\begin{document}

\portada{8}{Relación de potencia pico a promedio en sistemas OFDM}

\textbf{Script:} \texttt{practica8.py} (Python con NumPy/SciPy)

\section{Objetivo}

Analizar la relación de potencia pico a promedio (PAPR) de señales OFDM, calcular su valor
para distintos números de subportadoras y estudiar su distribución estadística mediante la
CCDF.

\section{Fundamento teórico}

La PAPR de una señal se define como
\begin{equation}
  \mathrm{PAPR} = \frac{\max_n |x[n]|^2}{\mathbb{E}\{|x[n]|^2\}},
  \qquad
  \mathrm{PAPR}_{\mathrm{dB}} = 10\log_{10}(\mathrm{PAPR}).
\end{equation}
En OFDM, $x[n]$ es la suma de $N$ fasores independientes. Por el teorema central del
límite, la amplitud tiende a una distribución gaussiana compleja, de modo que la
probabilidad de que muchas subportadoras se sumen en fase no es despreciable. La PAPR
crece aproximadamente como $\ln N$ y su valor máximo teórico es $N$ (10$\log_{10}N$ dB),
aunque en la práctica los valores observados son mucho menores.

Como la PAPR es aleatoria, se describe con la función de distribución complementaria
acumulada:
\begin{equation}
  \mathrm{CCDF}(\gamma) = P\{\mathrm{PAPR} > \gamma\},
\end{equation}
que indica la fracción de símbolos OFDM cuya PAPR supera un umbral $\gamma$.

Un amplificador de potencia debe operar con un margen (back-off) respecto de la
saturación para no recortar los picos. Ese margen reduce la eficiencia energética y es una
de las desventajas prácticas más importantes de OFDM.

\section{Desarrollo experimental}

\subsection{Parámetros}

\begin{tabularx}{\linewidth}{@{}*{2}{>{\raggedright\arraybackslash}X}@{}}
  \toprule
  Parámetro & Valor \\
  \midrule
  Subportadoras evaluadas & 16, 32, 64, 128 y 256 \\
  Realizaciones por configuración & 3000 \\
  Modulación & QPSK \\
  Distribución de la CCDF & 3000 símbolos por $N_{\mathrm{FFT}}$ \\
  Umbrales & 0 a 15 dB en pasos de 0.1 dB \\
  \bottomrule
\end{tabularx}

\subsection{Procedimiento}

\begin{enumerate}
  \item Para cada $N_{\mathrm{FFT}}$ se generan símbolos QPSK aleatorios y se aplica la
        IFFT con normalización de potencia unitaria.
  \item Se calcula la PAPR de cada símbolo y se promedian 3000 realizaciones.
  \item Se grafica un símbolo de 64 subportadoras para ver los picos en el tiempo y la
        potencia instantánea.
  \item Se construye la CCDF para 64, 128 y 256 subportadoras.
\end{enumerate}

\section{Código}

\begin{codigo}
def senal_ofdm(rng, nfft):
    x = np.fft.ifft(simbolo_qpsk(rng, nfft)) * np.sqrt(nfft)
    return x

def papr_db(x):
    p = np.abs(x) ** 2
    return float(10 * np.log10(np.max(p) / np.mean(p)))

ccdf = np.array([np.mean(papr_samples > u) for u in umbrales])
\end{codigo}

\section{Resultados}

\subsection{Potencia instantánea}

\figura{figs/potencia_instantanea.png}{Símbolo OFDM de 64 subportadoras y su potencia instantánea. La línea punteada es la potencia media.}

La potencia instantánea fluctúa alrededor de la media con picos que la superan varias
veces. En este símbolo la PAPR fue de 7.84 dB, es decir, el pico tuvo casi seis veces la
potencia media.

\subsection{PAPR en función del número de subportadoras}

\begin{tabular}{@{}cccc@{}}
  \toprule
  $N_{\mathrm{FFT}}$ & PAPR media (dB) & Desviación (dB) & PAPR máxima (dB) \\
  \midrule
  16  & 5.29 & 1.15 & 9.59 \\
  32  & 6.07 & 1.08 & 10.21 \\
  64  & 6.71 & 1.00 & 11.48 \\
  128 & 7.30 & 0.88 & 11.32 \\
  256 & 7.80 & 0.80 & 12.35 \\
  \bottomrule
\end{tabular}

\figura{figs/papr_vs_nfft.png}{PAPR media con barras de desviación estándar.}

La PAPR media crece de 5.29 a 7.80 dB al pasar de 16 a 256 subportadoras, un aumento de
2.51 dB para un factor 16 en $N$. El crecimiento es aproximadamente logarítmico, como
predice el modelo gaussiano. La dispersión alrededor de la media disminuye porque, con más
subportadoras, el promedio estadístico se estabiliza.

\subsection{CCDF}

\figura{figs/ccdf.png}{CCDF de la PAPR para 64, 128 y 256 subportadoras.}

Para una probabilidad de $10^{-2}$, la PAPR umbral es 9.4 dB con 64 subportadoras, 9.7 dB
con 128 y 10.1 dB con 256. Es decir, en el 1\,\% de los símbolos la PAPR supera esos
valores, y el amplificador debe reservar ese margen si no se aplica ninguna técnica de
reducción.

\section{Análisis de resultados}

\begin{itemize}
  \item \textbf{¿Por qué OFDM presenta valores elevados de PAPR?} Porque la señal es la suma
        coherente de muchas subportadoras. Cuando varias llegan en fase, sus amplitudes se
        suman y producen un pico mucho mayor que el valor eficaz.
  \item \textbf{¿Qué ocurre al incrementar el número de subportadoras?} La PAPR media
        aumenta de forma logarítmica y la CCDF se desplaza hacia la derecha: hay más
        ocasiones de alineación constructiva. En la simulación pasó de 5.29 a 7.80 dB.
  \item \textbf{¿Cómo afecta la PAPR a los amplificadores de potencia?} Obliga a operar con
        back-off para no recortar los picos. Un recorte produce distorsión no lineal,
        radiación fuera de banda y errores de demodulación.
  \item \textbf{¿Qué impacto tiene la PAPR sobre la eficiencia energética?} El amplificador
        trabaja lejos de su punto de máxima eficiencia; una PAPR de 8 dB puede reducir la
        eficiencia a la mitad o menos, lo que aumenta el consumo y el calor en la estación
        base y en el terminal.
\end{itemize}

\section{Cuestionario}

\begin{enumerate}
  \item \textbf{¿Qué es la PAPR?} La relación entre la potencia pico y la potencia promedio
        de una señal. Se expresa en dB y mide cuánto se alejan los picos del valor eficaz.
  \item \textbf{¿Por qué es un problema en OFDM?} Porque la suma de muchas subportadoras
        produce picos grandes con probabilidad no despreciable, lo que obliga a usar
        amplificadores con mucho margen y baja eficiencia.
  \item \textbf{¿Qué sistemas de comunicación utilizan OFDM?} LTE, 5G NR, Wi-Fi
        (802.11a/g/n/ac/ax), DVB-T, WiMAX, ADSL y comunicaciones por línea eléctrica.
  \item \textbf{¿Qué sucede si el amplificador opera en saturación?} Recorta los picos,
        genera distorsión no lineal, ensancha el espectro (interferencia fuera de banda) y
        aumenta la BER. Puede además dañar el amplificador por sobrecalentamiento.
  \item \textbf{¿Qué técnicas existen para reducir la PAPR?} El recorte (clipping) con
        filtrado, el mapeo selectivo (SLM), la transmisión de secuencias parciales (PTS),
        la reserva de tonos, la codificación de bloques y técnicas de precodificación como
        DFT-spreading (SC-FDMA).
\end{enumerate}

\section{Conclusiones}

Se midió la PAPR de señales OFDM QPSK para 16 a 256 subportadoras. La PAPR media creció de
5.29 a 7.80 dB y la CCDF mostró que el 1\,\% de los símbolos supera 9.4 dB con 64
subportadoras y 10.1 dB con 256. El resultado confirma el carácter probabilístico de los
picos y la necesidad de reservar margen en el amplificador o aplicar técnicas de reducción
de PAPR.

\end{document}
